\chapter{The other Bob testing charts for $q$ infinitesimal}
\label{component:pd2ecc6ac1891}
\noindent\textbf{Component }\path{pd2ecc6ac1891}\par
{\small\noindent Source: \path{R13/advances/c1_small/q_inf_remote_completion.tex}\par}
{\small\noindent SHA256: \path{d2ecc6ac1891279167b66f70a02259edbbec8f2f424b13387bbb0042e5a390b4}\par}
{\small\noindent Inventory use: QR1--QR3 and all homogeneous endpoint frames | C1\_\allowbreak{}BARE\_\allowbreak{}OTHER\_\allowbreak{}TERMINALS; precise subordinate uses retained below\par}
\medskip
\section{The other Bob testing charts for $q$ infinitesimal}
\label{pd2ecc6ac1891:sec:r7-c1-q-inf-remote}
\begin{theorem}
The original C1.1 source with $q=k^2\eta$ infinitesimal is tensor SOS
at every independent test. Together with the $q$-unit theorem this
closes the entire original C1.1 source.
\end{theorem}
The inherited $C,A$-large admission already handles $u$ not
infinitesimal. The preceding local theorem handles $u,v$
infinitesimal. Here are the two remaining $v$ charts. Every positive
opening rate and the literal permitted coefficient boundary is kept.

\paragraph{Ordinary $v$: the full normal residue is strict.}
Let $v$ be a positive bounded unit, use $Z=1$ on the negative-core
chart, and retain the roots of (QI1). Normalize $N$ by $H_0$ as
before, but now restore the positive packet $4h_b(vN)^2/t$.
Write
\[
 \tau=\operatorname{res}t,\qquad
 \rho=\operatorname{res}(t/e^2),\qquad
 a^2=\frac{\tau}{2(1+\tau)},\quad
 \gamma^2=\frac{\rho}{2(1+\tau)}.
\]
The original $\beta\ge0$ equation implies $0\le\rho\le1$;
$0\le\tau\le1$ follows from $M=R-t>0$ and $R\sim1$.
The entire residue on $(\widehat N,B,C)$ has diagonal entries
\[
 1+\frac{v^2}{1+\tau},\qquad 2(1+\tau v^2),\qquad 2,
\]
and polynomial cross coefficients
\[
 4av^2,\qquad -4\gamma v,\qquad -4\sqrt{\tau\rho}\,v.
\]
Its exact square decomposition is
\begin{align*}
 &2(C-\gamma v\widehat N-\sqrt{\tau\rho}\,vB)^2
       +\widehat N^2+2B^2\\
 &\qquad +(1-\rho)v^2
    \left(\frac{\widehat N}{\sqrt{1+\tau}}+
                         \sqrt{2\tau}\,B\right)^2.
 \tag{QR1}
\end{align*}
Thus it is strictly positive definite, also at $\rho=1$, the
possible $\beta=0$ endpoint, and at $\tau=0$ or $1$.
The other residue blocks are exactly
\[
                       2E_*^2+2(vh)^2+4D_vN_y^2.
\]
There are no nonzero leading cross entries with these blocks.

For completeness, the boundedness check uses the entire raw scalar
(QI2) and both full normal formulas (SC14), with (QI9) restored.
Terms containing $u h$ use $u h=uE_*+bC$; all their coefficients
are infinitesimal. The $N$--$B$ and $N$--$C$ coefficients with
nonzero residue are precisely those retained in (QR1). The harmonic
normal has infinitesimal coefficient on $h$ for bounded $v$:
after the opening cancellation its coefficient lies in
$(b^2,q,u^2)$ with bounded arrays on this chart. Its other
coefficients are infinitesimal as well.
All remaining source-constant terms and weak monomials enter the
same strict Gram through the exact bounded pure-testing identities
\[
                   h=(vh)/v,\qquad
                   1=((vh)-vE_*)/v^2.
\]
For example $A=ch$ and $D=qg$ are infinitesimal-coefficient
combinations of these two roots. Consequently the actual full Gram
has the positive residue just displayed. Restore all retained
packets. The same sign, $Z=0$, and complex-state joins as in the
local theorem apply; no restriction on the size of the normalized
state was used.

\paragraph{Unbounded $v$: a five-root endpoint frame.}
Put $w=1/v$, so $w$ is infinitesimal, and multiply the original
parent by $w^2$. The exact source formulas give positive real
constants $\mu,d_0$ with
\[
 C/d\ge\mu/b^2,\qquad A/d\ge\mu/b^2,\qquad
 D/d\ge d_0 b^2.
\]
For instance sufficiently small fixed constants below $1/4$ work,
since $b^2C/d,b^2A/d$ are uniformly bounded below and
$D/(db^2)\sim1/2$. Retain all excesses and the entire $m,\alpha,
\beta$ packets. No assumption on $e$ is needed.
Use the five roots
\[
              P=uW/b,\quad Q=Y/b,\quad X,\quad J_b'=bZ,
                     \quad V_b=uwZ.
\]
The homogenized signed core is exactly
\begin{align*}
 K={}&V_b^2-2buV_bQ-2wPJ_b'+2wV_bX+b^2w^2P^2
       +2buwPX-2b^2w^2PQ\\
 &+(D_u+w^2)(X^2+b^2Q^2)-2buwXQ.
 \tag{QR2}
\end{align*}
The three selected event packets are
\[
 4\mu D_w(P-Q)^2+4\mu D_uD_wQ^2,
 \qquad
 4d_0(b^2wP+buX-b^2wQ-J_b')^2.
 \tag{QR3}
\]
Combining these with $2K$, the bounded Gram has strictly positive
residue
\[
 4\mu\{(P-Q)^2+Q^2\}+4d_0(J_b')^2+2X^2+2V_b^2.
\]
All other displayed coefficients have zero residue. This is a full
state-quadratic certificate, not a restriction to a contact graph.
It includes the homogeneous endpoint $w=0$ directly and uses no
inverse involving independent source and testing variables.

The infinitesimal, unit, and unbounded $v$ charts exhaust all tests
remaining after the inherited $C,A$ admission. Hence the entire
$q$-infinitesimal C1.1 source, and consequently C1.1 itself, is closed.
